%! AP Statistics — Unit 4, Lesson 4.8 — Guided Notes
\documentclass[10pt]{article}
\usepackage{apstats-article}
\usepackage{apstats-boxes}
\newcommand{\TallMath}[1]{$\displaystyle #1\rule[-1.4em]{0pt}{3.2em}$}

\begin{document}

\pageheader{Unit 4, Lesson 4.8}{Guided Notes}
\namedateperiod

\vspace{0.12in}

\begin{objectivebox}
\small By the end of this lesson, I will be able to\dots\\[4pt]
\writeline\\[1pt]
\writeline\\[1pt]
\writeline
\end{objectivebox}

\vspace{0.1in}

\begin{vocabbox}
\termblanklong{Expected value ($\mu_X$)}
\termblanklong{Variance ($\sigma^2_X$)}
\termblanklong{Standard deviation ($\sigma_X$)}
\termblanklong{Long-run average}
\end{vocabbox}

\vspace{0.1in}

\begin{hookbox}
\small
A carnival game costs \$2 to play. Prize distribution:
\begin{center}
\renewcommand{\arraystretch}{1.3}
\begin{tabular}{@{}ccc@{}}
Prize ($x$) & \$0 & \$5 \\
\hline
Probability & 0.70 & 0.30 \\
\end{tabular}
\end{center}
If you play 100 times, what do you expect to win in total? \blank{4cm}\\
Is playing worth the \$2 entry fee? \blank{5cm}
\end{hookbox}

\vspace{0.1in}

\begin{notesbox}{Mean (Expected Value) of a Random Variable}
\small

The \textbf{expected value} of $X$ is a \blank{3cm} average of all possible values,
weighted by their probabilities:

\vspace{8pt}
\begin{center}
\TallMath{\mu_X = \sum x \cdot P(X=x) = x_1 \cdot P(X=x_1) + x_2 \cdot P(X=x_2) + \cdots}
\end{center}

\vspace{8pt}
\textbf{Computation Table Example:} $X$ = prize from carnival game.

\vspace{4pt}
\renewcommand{\arraystretch}{1.8}
\begin{tabularx}{\linewidth}{@{} c c X @{}}
$x$ & $P(X=x)$ & $x \cdot P(X=x)$ \\
\hline
\$0 & 0.70 & \blank{4cm} \\
\$5 & 0.30 & \blank{4cm} \\
\end{tabularx}

\vspace{6pt}
$\mu_X = \sum x \cdot P(X=x) = $ \blank{4cm}

\vspace{6pt}
\textbf{Interpret} $\mu_X$ in context: \blank{3cm} means that in the \blank{4cm},\\
the average prize per play is approximately \blank{2.5cm}.

\end{notesbox}

\vspace{0.1in}

\begin{notesbox}{Standard Deviation of a Random Variable}
\small

The \textbf{variance} measures how spread out the values of $X$ are:

\vspace{6pt}
\begin{center}
\TallMath{\sigma^2_X = \sum (x - \mu_X)^2 \cdot P(X=x)}
\end{center}

\vspace{6pt}
The \textbf{standard deviation} is $\sigma_X = $ \blank{3cm}.

\vspace{8pt}
\textbf{Full Computation Table:} $X$ = number of cars per household ($\mu_X = 1.55$).

\vspace{4pt}
\renewcommand{\arraystretch}{2.0}
\begin{tabularx}{\linewidth}{@{} c c c c X @{}}
$x$ & $P(X=x)$ & $x \cdot P(X=x)$ & $(x - \mu_X)^2$ & $(x-\mu_X)^2 \cdot P(X=x)$ \\
\hline
0 & 0.10 & \blank{1.5cm} & \blank{1.8cm} & \blank{2.5cm} \\
1 & 0.40 & \blank{1.5cm} & \blank{1.8cm} & \blank{2.5cm} \\
2 & 0.35 & \blank{1.5cm} & \blank{1.8cm} & \blank{2.5cm} \\
3 & 0.15 & \blank{1.5cm} & \blank{1.8cm} & \blank{2.5cm} \\
\hline
\textbf{Sum} & & $\mu_X =$ \blank{1cm} & & $\sigma^2_X =$ \blank{1.5cm} \\
\end{tabularx}

\vspace{6pt}
$\sigma_X = \sqrt{\sigma^2_X} = $ \blank{3cm}

\vspace{6pt}
\textbf{Interpret} $\sigma_X$ in context:\\
\writeline

\end{notesbox}

\vspace{0.1in}

\begin{practicebox}
\small
$X$ = number of goals scored per soccer game.
\begin{center}
\renewcommand{\arraystretch}{1.4}
\begin{tabular}{@{}cccc@{}}
$x$ & 0 & 1 & 2 \\
\hline
$P(X=x)$ & 0.30 & 0.50 & 0.20 \\
\end{tabular}
\end{center}

\renewcommand{\arraystretch}{2.0}
\begin{tabularx}{\linewidth}{@{} c c c c X @{}}
$x$ & $P(X=x)$ & $x \cdot P(X=x)$ & $(x - \mu_X)^2$ & $(x-\mu_X)^2 \cdot P(X=x)$ \\
\hline
0 & 0.30 & \blank{1.5cm} & \blank{1.8cm} & \blank{2.5cm} \\
1 & 0.50 & \blank{1.5cm} & \blank{1.8cm} & \blank{2.5cm} \\
2 & 0.20 & \blank{1.5cm} & \blank{1.8cm} & \blank{2.5cm} \\
\hline
\textbf{Sum} & & $\mu_X =$ \blank{1cm} & & $\sigma^2_X =$ \blank{1.5cm} \\
\end{tabularx}

\vspace{6pt}
$\sigma_X = $ \blank{3cm}

\vspace{6pt}
Interpret $\mu_X$ in context: \writeline

\vspace{2pt}
Interpret $\sigma_X$ in context: \writeline
\end{practicebox}

\end{document}
